\section{Two checks on a conditioning formula}\label{sec:conditioning}

The following examples concern equation (62) on page 38 of Bordenave's arXiv version 4 \cite{Bordenave}. They refute that intermediate conditional assertion under its displayed conditioning. They are not counterexamples to Proposition 28, which Section~\ref{sec:mixed} proves by a different argument.

For requested cells $y_t=(i_t,j_t)$, the source conditions on
\[
\mathcal F_t=\sigma\bigl(\sigma(i_s),\sigma^{-1}(j_s):s\ne t\bigr)
\]
and asserts a probability of the form
\[
\PP(M_{i_tj_t}=1\mid\mathcal F_t)
=\frac{\1_{\Omega_t^c}}{n_t}\le\frac1{n-a+1},
\qquad
n_t=n-|\{\sigma(i_s):s\ne t\}|.
\]
The event $\Omega_t$ holds if some $s\ne t$ satisfies $\sigma(i_s)=j_t$ or $\sigma^{-1}(j_s)=i_t$.

\begin{example}[A shared row reveals the cell]
Let $n=4$, $y_1=(1,1)$ and $y_2=(1,2)$. Then $a=k=2\le\sqrt n$, and $\mathcal F_1$ includes $\sigma(1)$. On the event $\sigma(1)=1$, of probability $1/4$, the actual conditional probability of $M_{11}=1$ is one. The displayed event $\Omega_1$ also holds, because $\sigma(i_2)=j_1$. The source formula therefore gives zero. Its claimed upper bound $1/3$ fails as well.
\end{example}

\begin{example}[Images and preimages remove two arrows]
Let $n=4$, $y_1=(1,1)$ and $y_2=(2,2)$. Condition on
\[
\sigma(2)=3,\qquad \sigma^{-1}(2)=4.
\]
This event has probability $1/12$. It exposes two arrows, $2\mapsto3$ and $4\mapsto2$. The remaining rows are $1,3$ and the remaining columns are $1,4$. Thus
\[
\PP(M_{11}=1\mid\mathcal F_1)=\frac12.
\]
Neither clause of $\Omega_1$ holds, but the displayed $n_t$ is $4-1=3$. The source formula gives $1/3$, and the proposed bound $1/(n-a+1)$ also gives $1/3$.
\end{example}

The second example shows that merely requiring the tested cell to be unrevealed does not repair the formula. One must count every removed row and column under the actual conditioning. Section~\ref{sec:mixed} avoids the issue by exposing only a compatible prescribed set. Its remaining dimension is exactly $n-r$.
